\section{Prejudices of Philosophers}

\begin{quotex}
Having kept a sharp eye on philosophers, and having read between their lines long enough, I now say to myself that the greater part of conscious thinking must be counted among the instinctive functions, and it is so even in the case of philosophical thinking; one has here to learn anew, as one learned anew about heredity and ``innateness.'' As little as the act of birth comes into consideration in the whole process and procedure of heredity, just as little is ``being-conscious'' OPPOSED to the instinctive in any decisive sense; the greater part of the conscious thinking of a philosopher is secretly influenced by his instincts, and forced into definite channels. And behind all logic and its seeming sovereignty of movement, there are valuations, or to speak more plainly, physiological demands, for the maintenance of a definite mode of life.
\flright{\textsc{Friedrich Nietzsche}}
\end{quotex}


\paragraph{Multiple Selves}


\begin{quotex}
It has gradually become clear to me what every great philosophy up till now has consisted of---namely, the confession of its originator, and a species of involuntary and unconscious auto-biography ... our body is but a social structure composed of many souls.
\flright{\textsc{Friedrich Nietzsche}}
\end{quotex}


\textbf{Jordan Peterson} spoke for 45 minutes on the first chapter of \emph{Beyond Good and Evil}\footnote{\url{https://www.youtube.com/watch?v=MCOw0eJ84d8}}. Philosophers, and everyone else for that matter, believe that they are following an impulse to knowledge. Nietzsche challenges that with the assertion that unconscious impulses---perhaps genii, demons, or kobolds---are the real sources of their ``philosophy''. But long before Nietzsche, or Jung, or even Peterson, that was the teaching of the esoteric schools. Nietzsche goes further:

\begin{quotex}
With regard to the superstitions of logicians, I shall never tire of emphasizing a small, terse fact, which is unwillingly recognized by these credulous minds---namely, that a thought comes when ``it'' wishes, and not when ``I'' wish; so that it is a PERVERSION of the facts of the case to say that the subject ``I'' is the condition of the predicate ``think.''
\flright{\textsc{Friedrich Nietzsche}}
\end{quotex}


Again, the esoteric schools developed exercises to locate the sources of the spontaneous thoughts that arise in the mind. For example, Valentin Tomberg explains the foundation of \textbf{Patanjali}'s \emph{Yoga Sutras}, a summary of knowledge from several millennia ago:

\begin{quotex}
The ``oscillations'' (\emph{vritti}) of the ``mental substance'' (\emph{citta}) take place automatically. This automatism in the movements of thought and imagination is the opposite of concentration.
\end{quotex}


Nietzsche is close with the idea of genii. Our thoughts may arise, for example, from lower luciferic forces, bodily needs, the collective and personal unconscious, or even from transcendent sources. The higher human soul functions are the intellect and the will. Since Nietzsche rejects the ideal of the Will to Truth, he centres his philosophy in the Will rather than the mind:

\begin{quotex}
Psychologists should bethink themselves before putting down the instinct of self-preservation as the cardinal instinct of an organic being. A living thing seeks above all to DISCHARGE its strength---life itself is WILL TO POWER.
\end{quotex}


But what exactly does the Will to Power will? Nietzsche criticizes the notion of an Ego that thinks:

\begin{quotex}
What gives me the right to speak of an `ego', and even of an `ego' as cause, and finally of an `ego' as cause of thought?
\end{quotex}


That is why the esoteric schools teach the very opposite, viz., inner silence rather than thinking. Tomberg explains:

\begin{quotex}
The ``oscillations of the mental substance'' will never be able to be reduced to silence if the will itself does not infuse them with its silence. It is the silenced will which effects the silence of thought and imagination in concentration. This is why the great ascetics are also the great masters of concentration. All this is obvious and stands to reason.
\end{quotex}


Concentration is actually the will to power:

\begin{quotex}
True concentration is a free act in light and in peace. It presupposes a disinterested and detached will. For it is the condition of the will which is the determining and decisive factor in concentration.
\end{quotex}


Nietzsche repeated Epicurus' critique of the Platonists as ``actors''; we would say today ``larpers''. Ironically, the same could be said of Nietzscheans today: if they are not self-critical, they are larpers.

\paragraph{The Hard Problem of Consciousness}


In an interview with \textbf{Roger Penrose}\footnote{\url{https://www.youtube.com/watch?v=GEw0ePZUMHA}}, Joe Rogan claims to have read Penrose's books.

Even if they are Penrose's more accessible books, that is still surprising. One of Penrose's pet projects is to make consciousness depend on quantum gravity, albeit with no proof whatsoever. There is still no convincing scientific or positivist theory of consciousness; that is why it is considered to be a ``hard problem''.

But if science is based on the empirical evidence of the senses, then consciousness can never be explained scientifically. One can observe one's own consciousness, but not that of another. Another's consciousness has to be presumed indirectly based on reports and behaviour. That is the point of the Turing Machine\footnote{\url{https://gornahoor.net/?p=8390}}. It would make more sense to accept that consciousness is fundamental to the universe and not reducible to anything else, particularly anything material.

\paragraph{Animal Intelligence}


\begin{quotex}
But the sensual or animal man perceiveth not these things that are of the Spirit of God; for it is foolishness to him, and he cannot understand, because it is spiritually examined. But the spiritual man judgeth all things; and he himself is judged of no man.
\flright{\textsc{I Corinthians 2: 14-15}}
\end{quotex}


The next topic discussed is the idea of animal intelligence. The two seemed quite impressed at the intelligence of animals, which they deemed comparable to humans. The two are unaware that Traditional thinking has always known that humans and animals shared a sensitive soul, so there should be little surprise to see complex behaviour in the beasts. After all, the World was created by Logos, so intelligence is foundational. We address this topic a few year ago in \textit{Scientific and Political Myths}\footnote{\url{https://gornahoor.net/?p=7827}}:

\begin{quotex}
There has been a story circulating in the news about a dog that allegedly gets on a bus alone in order to get to a dog park. People find it fascinating because it would seem to indicate some measure of the dog's human-like intelligence. On the contrary, the real lesson is how little intelligence it takes humans to get through their day. So if you are getting on the bus at the same stop and getting off at another, day after day, you have demonstrated the intelligence of a dog.
\end{quotex}


The issue is that the intelligence of animals does not choose the ends but only the means. The fundamental ends of animal life are Fear, Sex, Hunger, i.e., repulsion, attraction, and assimilation. Hence, a pigeon can be taught to solve a sequence of seemingly complex puzzles in order to be rewarded with a food pellet; this it will do over and over again. However, it won't be as motivated to learn the Fundamental Theorem of Algebra\footnote{\url{https://en.wikipedia.org/wiki/Fundamental\_theorem\_of\_algebra}}.

Then many animal species will engage in complex mating rituals, even risky ones due to competition between males. The threat of danger is also associated with intelligence, including various ruses. Group behaviour needs to follow mathematical Game Theory, since the World is subject to mathematical laws. Otherwise, the species does not survive.

The point is that animal intelligence is purely instrumental: it is subordinated to the life forces and feelings of the animal. It has no value in itself. Animal intelligence is fixed and repeatable. For example, a bee in a beehive will never rebel against the social structure of the hive by fomenting a revolution, unlike human societies. Bees always act in the same way.

\paragraph{Human Intelligence}


Aristotle famously defined man as a rational animal. As was shown above, man retains the animal intelligence. To that, however, is added the intellectual soul with the power of rationality. What exactly is it? Human rationality includes the ability to choose one's ends, so that man is not limited to the primary motives of fear, sex, and hunger.

That ability to choose turns the world into a moral world. What is automatic in animals becomes a choice for humans. For example, God commands Noah to ``be fruitful and multiply.'' Animals need no such command, since they will naturally mate (other than perhaps in some extreme circumstances). Man also has the choice to ignore that command; fertility rates in most developed countries are below replacement level.

Hence, we cannot model ourselves on animal behaviour, as many would like to do. Man must choose ends intelligently, not based solely on emotions, unconscious drivers, or other lower psychic functions.

\paragraph{Psychoactive drugs}


If the intellectual soul is what distinguishes humans from animals, then it is inadvisable to use artificial means to suppress its proper functioning. This includes recreational drugs especially, whether the purpose is to dull consciousness or to accumulate experiences. Unfortunately, this also includes prescription medications, although it is better to take them than to suffer. For some conditions they are mandatory.

Nevertheless, the prevalence of drugs in the population has had an adverse effect on social discourse.


\flrightit{Posted on 2020-01-16 by Cologero}

\begin{center}* * *\end{center}

\begin{footnotesize}
\begin{sffamily}
\texttt{Anonymous on 2020-01-17 at 16:59 said:}

Certainly better thinkers could've been included to demonstrate examples than Joe Rogan and Jordan Peterson. I was quite disappointed with this article.

\hfill

\texttt{Cologero on 2020-01-18 at 21:02 said:}

The classical style represents essentially such tranquility, simplification, abbreviation, concentration --- in the classical type the supreme feeling of power is concentrated. Slow to react: a tremendous consciousness: no feeling of struggle. 
\flright{\emph{The Will to Power} \#799}


\hfill

\texttt{Cologero on 2020-01-18 at 23:14 said:}

Yes, and I am often quite disappointed in the comments I get.

\hfill

\texttt{Max on 2020-01-26 at 11:04 said:}

The same could be said of gravity as for consciousness: ``There is still no convincing scientific or positivist theory of gravity.'' Quantum gravity implies gravity-particles, or in other words, old tired atomism resuscitated again. This was a major battle-line between philosophers of antiquity. The theoretical struggle for quantification, or what Guénon calls ``the reign of quantity'', has to be located within a larger framework of will to power, even if it really represents an inversion through the rejection of agency.

Scientists attempts to reduce gravity to something physical, but that is a dead end. Gravity describes the phenomenon of bodies ``falling'' towards one another, that is to say, it indicates an absence rather than a presence, or the anti-thesis of space. The experience of gravity is free fall, which does not feel like anything at all, since sensation arises upon impact against another free-falling body, resisting it.

Platonists who are self-conscious of the will-to-power of a philosopher-king does not seem as self-deluded as actors merely playing out someone else's drama. Nietzsche speaks of dogmatic Platonism, although that is only half the story, since the forgotten half is made up of skeptic Platonism.

\hfill

\end{sffamily}
\end{footnotesize}

